\documentclass[11pt]{article}
\usepackage[margin=1in]{geometry}
\usepackage{amsmath,amssymb,bm}
\title{Conduction inside the circuit}
\author{}
\date{}
\begin{document}
\maketitle

\section{Problem Setup}
The conductor is one closed loop, divided into wire, one battery, and
possibly several resistor regions. Let $\rho$ be the resistivity in each
region and $\gamma=1/\rho$ its conductivity. The electrostatic field is
$\bm E=-\nabla V$. The battery also pushes charge by a mechanism represented
by a driving field $\bm F$. The current density is
\begin{equation}
 \bm J=\gamma(\bm E+\bm F)
       =\gamma(-\nabla V+\bm F).
 \label{eq:current}
\end{equation}
If the battery has EMF $\mathcal E$, length $\ell_b$, and unit driving
direction $\bm d_b$, the model sets
\begin{equation}
 \bm F=\begin{cases}
       (\mathcal E/\ell_b)\bm d_b,&\text{inside the battery},\\
       0,&\text{in the wire and resistors}.
       \end{cases}
\end{equation}
The EMF is the total work supplied per unit charge across the battery;
$\bm F$ is that drive per unit length. It is separate from the electric
field $\bm E$ made by charges.

In steady state $\nabla\cdot\bm J=0$: charge does not build up inside
the conductor. Current cannot cross either conductor--air surface, so
$\bm J\cdot\bm n=0$ there. At a material junction, the potential and
normal current are continuous. No terminal voltage is prescribed,
because the conductor forms a closed loop.

\section{Turning conservation into equations}
The shared derivation in \texttt{fem.tex} starts with
$\bm a=c(\nabla u-\bm f)$ and $\nabla\cdot\bm a=0$. Here choose
$u=V$, $c=\gamma$, and $\bm f=\bm F$, so $\bm a=-\bm J$.
The boundary term from integration by parts vanishes because
$\bm J\cdot\bm n=0$ on both conductor surfaces. Hence, for each
allowed test function $w$,
\begin{equation}
 \int_{\Omega_c}\gamma\nabla V\cdot\nabla w\,dA
 =\int_{\Omega_c}\gamma\bm F\cdot\nabla w\,dA.
 \label{eq:weak}
\end{equation}
Triangles on both sides of a material junction share a common edge.
Thus the finite-element potential is continuous there, and
equation~\eqref{eq:weak} makes the normal currents balance.

\section{Choosing a potential reference}
If $v$ solves the conductor equations, then $v+C$ also solves them:
only its gradient appears. This is why the assembled matrix has
a redundant degrees of freedom.
The solver fixes one conductor vertex at zero
and solves the remaining free-vertex equations. It then chooses a more
useful reference by subtracting the mean potential over the conductor.
Let $\mathcal T_c$ be the set of conductor triangles. Then
\begin{equation}
 \overline V=
 \frac{\displaystyle\sum_{h\in\mathcal T_c}
       A_h(V_{h,0}+V_{h,1}+V_{h,2})/3}
      {\displaystyle\sum_{h\in\mathcal T_c}A_h},
 \qquad V\leftarrow V-\overline V.
\end{equation}
For a linear function on a triangle, the average of its three vertex
values is its area average. The final shift changes no electric field
or current.

\section{Fields used by the app}
After solving for the potential, the field and current in each
conductor triangle are
\begin{equation}
 \bm E_h=-\sum_{i=0}^2 V_{h,i}\nabla N_{h,i},\qquad
 \bm J_h=\frac{\bm E_h+\bm F_h}{\rho_h}.
\end{equation}
Both vectors are constant within a triangle. The app samples these
vectors to draw electric-field and current arrows. Air current is zero;
the electric field in air is calculated in \texttt{exterior.tex}.

\paragraph{Code:} \texttt{calculations/conduction.py}.
\end{document}
